\documentclass[11pt]{amsart}
\usepackage[T1]{fontenc}
\usepackage{lmodern,microtype,amsmath,amssymb,amsthm,mathtools,booktabs}
\usepackage[margin=1.1in]{geometry}
\usepackage[hidelinks]{hyperref}
\usepackage{xurl}

\newtheorem{theorem}{Theorem}
\newtheorem{proposition}[theorem]{Proposition}
\theoremstyle{remark}
\newtheorem{remark}[theorem]{Remark}

\newcommand{\dbar}{\bar d}
\newcommand{\TV}{\mathrm{TV}}

\hypersetup{pdftitle={The pairwise distance of Ising Glauber dynamics is not monotone in the couplings},pdfauthor={Yitzchak Shmalo}}

\title[Distance of Ising Glauber dynamics and couplings]{The pairwise distance of Ising Glauber dynamics is not monotone in the couplings}
\author{Yitzchak Shmalo}
\address{Einstein Institute of Mathematics, The Hebrew University of Jerusalem, Jerusalem, Israel}
\email{yitzchak.shmalo@gmail.com}
\date{November 2026}
\subjclass[2020]{60J10, 82C20, 60K35}
\keywords{Glauber dynamics, Ising model, mixing, total variation distance, monotonicity}

\begin{document}
\begin{abstract}
For the Glauber dynamics of a ferromagnetic Ising model on a finite graph, Levin and Peres ask whether, for a fixed graph and a fixed time $t$, the largest total variation distance between the laws at time $t$ from two starting configurations increases in each interaction strength. It does not. On the path with four vertices and couplings $(2,x,4)$, raising $x$ from $1$ to $2$ lowers this distance at every time $t$ from $7$ to $699$, and raises it at $t=700$; these inequalities are certified by exact interval arithmetic. The same couplings give a decrease for the lazy chain and in continuous time. The spectral gap behaves as monotonicity of the gap requires: it shrinks when $x$ grows, and the prefactor of the leading term decreases enough to outweigh it until $t=700$. A second example lives on the $5$-cycle, where the spectral gap is known to be monotone.
\end{abstract}
\maketitle

\section{Introduction}\label{sec:intro}

Let $G=(V,E)$ be a finite graph and let $J=(J_{uv})_{uv\in E}$ be positive couplings. The Ising measure is
\[
\mu_J(\sigma)=\frac1{Z_J}\exp\Bigl(\sum_{uv\in E}J_{uv}\,\sigma_u\sigma_v\Bigr),\qquad \sigma\in\{-1,1\}^V,
\]
and its Glauber (heat-bath) dynamics picks a vertex $v$ uniformly at random and resamples $\sigma_v$ from its conditional law given the other spins:
\begin{equation}\label{eq:cond}
\mu_J(\sigma_v=+1\mid\sigma_{V\setminus v})=\frac{1}{1+e^{-2h_v(\sigma)}},\qquad h_v(\sigma)=\sum_{u\colon uv\in E}J_{uv}\sigma_u .
\end{equation}
Write $P_J$ for the transition matrix and
\[
\dbar_J(t)=\max_{\sigma,\tau}\bigl\|P_J^t(\sigma,\cdot)-P_J^t(\tau,\cdot)\bigr\|_{\TV}.
\]

Let $1=\lambda_1(J)>\lambda_2(J)\ge\cdots$ be the eigenvalues of $P_J$, so that $1-\lambda_2(J)$ is the spectral gap. Levin and Peres~\cite[Section~26.1, Question~2]{LP} ask whether the spectral gap is nonincreasing in each coupling $J_{uv}$; Nacu~\cite{Nacu} proved this for the cycle with arbitrary couplings. They then ask for more: whether, for a fixed graph and a fixed $t$, the distance $\dbar_J(t)$ is nondecreasing in each coupling. As they note, a positive answer would settle the spectral-gap question, because $P_J$ has no negative eigenvalue (Section~\ref{sec:mech}) and so $\dbar_J(t)^{1/t}\to\lambda_2(J)$ as $t\to\infty$~\cite[Corollary~12.7 and Lemma~4.10]{LP}. Shirai~\cite{Shirai} showed that on the path with four vertices, with couplings $(2s,x,s)$ and $s$ large, the expectation of a single spin at time $3$, started from all plus, is not monotone in $x$. By spin-flip symmetry this expectation is a lower bound for $\dbar_J(t)$, and a lower bound can decrease while $\dbar_J(t)$ increases, so the question for $\dbar_J(t)$ remains.

On the same path the distance itself decreases. When $\lambda_2(J)$ is a simple eigenvalue,
\[
\dbar_J(t)=C_J\,\lambda_2(J)^t\bigl(1+o(1)\bigr),\qquad t\to\infty,
\]
where $C_J$ is the largest total variation distance between two rows of the spectral projector onto the eigenspace of $\lambda_2(J)$. For the couplings $(2,x,4)$, raising $x$ from $1$ to $2$ raises $\lambda_2$ from $0.999834342$ to $0.999837852$ and lowers $C$ from $1.003641$ to $1.001180$. The smaller prefactor wins until $(\lambda_2(2)/\lambda_2(1))^t$ catches up with $C_1/C_2$, near $t=700$.

\begin{theorem}\label{thm:main}
Let $G$ be the path $0-1-2-3$ with $J_{01}=2$, $J_{12}=x$, $J_{23}=4$, and write $\dbar_x$ for the corresponding distance. Then
\[
\dbar_2(t)<\dbar_1(t)\quad\text{for every integer } 7\le t\le 699,\qquad \dbar_2(700)>\dbar_1(700).
\]
\end{theorem}

The proof is the interval computation of Section~\ref{sec:cert}. Table~\ref{tab:main} gives the values at both ends of the window and at three times inside it.

\begin{table}[tbp]
\centering
\begin{tabular}{rlll}
\toprule
$t$ & $\dbar_1(t)$ & $\dbar_2(t)$ & $\dbar_2(t)-\dbar_1(t)$\\
\midrule
6 & $0.99959217877695345181$ & $0.99959672151288890647$ & $\phantom{-}4.5427359354547\times10^{-6}$\\
7 & $0.99948922026833168240$ & $0.99948446850266242769$ & $-4.7517656692547\times10^{-6}$\\
30 & $0.99679585729479756393$ & $0.99623502850285614656$ & $-5.6082879194142\times10^{-4}$\\
144 & $0.97979763674586148394$ & $0.97807158399996023267$ & $-1.7260527459013\times10^{-3}$\\
699 & $0.89389222242739674918$ & $0.89389046455434841506$ & $-1.7578730483341\times10^{-6}$\\
700 & $0.89374414211522642175$ & $0.89374552169994777552$ & $\phantom{-}1.3795847213538\times10^{-6}$\\
\bottomrule
\end{tabular}
\caption{The path with couplings $(2,x,4)$, discrete-time heat-bath dynamics, at the two ends of the window and at three times inside it. Entries are rounded to the digits shown; the enclosures of Section~\ref{sec:cert} agree with them to more digits than are printed.}\label{tab:main}
\end{table}

For $2\le t\le6$ the order is also reversed, $\dbar_2(t)>\dbar_1(t)$, and at $t=1$ both distances equal $1$. In double precision the reversed order persists for every $t$ from $700$ to $3000$. At $t=7$, $30$ and $144$ the maximum in $\dbar_x(t)$ is attained at the pair (all plus, all minus).

\begin{remark}[Lazy chain and continuous time]\label{rem:lazy}
For the lazy chain $(I+P_x)/2$ the same certificate gives decreases of $1.5859546851\times10^{-4}$ at $t=30$ and $1.3851503307\times10^{-3}$ at $t=144$; numerically the window is $12\le t\le1399$. In continuous time, with each vertex updated at rate one (generator $4(P_x-I)$), the distance at time $t$ drops by $2.9\times10^{-4}$, $8.1\times10^{-4}$, $1.47\times10^{-3}$, $1.69\times10^{-3}$, $1.71\times10^{-3}$ and $1.53\times10^{-3}$ at $t=5$, $10$, $20$, $30$, $40$, $60$ when $x$ goes from $1$ to $2$.
\end{remark}

\begin{remark}[The cycle]\label{rem:cycle}
On the $5$-cycle with couplings $(x,2,\tfrac12,4,2)$ around the cycle, raising $x$ from $\tfrac12$ to $1$ lowers $\dbar(3)$ from $0.99997418786773027079$ to $0.99997214160667527597$, computed at $80$ significant digits over all $496$ pairs of states. So the fixed-time monotonicity fails even on the graph where Nacu proved that the spectral gap is monotone.
\end{remark}

\section{The mechanism}\label{sec:mech}

The heat-bath kernel is an average of the $|V|$ single-site updates, each an orthogonal projection in $L^2(\mu_J)$, so $P_J$ is self-adjoint with spectrum in $[0,1]$. For the path of Theorem~\ref{thm:main} the second eigenvalue is
\[
\lambda_2(1)=0.99983434204549494558,\qquad \lambda_2(2)=0.99983785166063618186,
\]
so the spectral gap shrinks as $x$ grows, as a monotone spectral gap requires. Both are simple eigenvalues, and the third largest eigenvalue is $0.97983$ for $x=1$ and $0.92200$ for $x=2$. The asymptotic formula of the introduction therefore applies, with
\[
C_1=1.00364104987661629280,\qquad C_2=1.00117954021809281471.
\]
Stronger coupling slows the exponential rate but lowers the prefactor. The two leading terms are equal when $t=\log(C_1/C_2)/\log(\lambda_2(2)/\lambda_2(1))=699.56$, and the distances themselves change order between $t=699$ and $t=700$ (Theorem~\ref{thm:main}). In particular the example says nothing against the spectral-gap question.

\section{The certificate}\label{sec:cert}

All couplings in Theorem~\ref{thm:main} are integers, so every local field $h_v$ is an integer and every entry of $P_x$ is a rational function of $e^{-2}$. We enclose $e^{-2}$ between two consecutive partial sums of its alternating Taylor series (degrees $200$ and $201$) and carry out all further arithmetic on intervals whose endpoints are integers in units of $10^{-100}$, rounding every product and division outward. This gives interval matrices containing $P_1$ and $P_2$, and multiplying by them one factor at a time gives interval matrices containing $P_x^t$ for $t=1,\dots,700$. For each pair of rows, a lower bound for the total variation distance is half the sum, over the $16$ columns, of the distance from $0$ to the interval of the entrywise difference, and an upper bound is half the sum of the largest absolute values in those intervals. The maximum over the $120$ pairs gives an interval for $\dbar_x(t)$ of width below $2\times10^{-96}$. For every $t$ from $2$ to $700$ the intervals for $x=1$ and $x=2$ are disjoint, in the order stated in Theorem~\ref{thm:main} and after it. No floating-point operation enters. The same computation with $(I+P_x)/2$ in place of $P_x$ gives the lazy-chain values in Remark~\ref{rem:lazy}.

The eigenvalues of Section~\ref{sec:mech}, the windows quoted in double precision and the values in the remarks other than the lazy-chain ones are computed numerically, in multiple precision ($50$ to $120$ significant digits) where digits are shown.

\section*{Acknowledgements and funding}
This work was done by AI agents connected to Math Brain~\cite{MathBrainRepo} and working inside the Brain framework~\cite{BrainFramework}, a shared, persistent memory in which the agents record claims, proofs, failed attempts and methods, retrieve one another's results, coordinate who works on what, and audit each other's work. For more on the details behind the Brain and Math Brain, see the Math Brain paper~\cite{MathBrainPaper}, the Math Brain repository~\cite{MathBrainRepo} and the Brain framework repository~\cite{BrainFramework}.

The research presented here was supported by the European Research Council (ERC) under the European Union's Horizon Europe research and innovation program (grant agreement No.~101041711) and by the Simons Foundation as part of the Collaboration on the Mathematical and Scientific Foundations of Deep Learning. It was also supported by Heights Labs, by Convex Nexus Capital, by the Israel Science Foundation (grant number 2258/19), and by the Israel Science Foundation (ISF Grant 4101/25). Computational resources were provided in part by Google Cloud through its education credits program (credit 530550103).

\section*{Declaration of competing interest}
The author holds equity in Heights Labs and in Convex Nexus Capital and has a management role at Convex Nexus Capital. These entities are included in the funding declaration above.

\section*{Declaration of generative AI and AI-assisted technologies}
In November 2026 Claude Opus~5.5 (Anthropic), working in Claude Code, found the examples, wrote and ran the numerical scripts, and wrote the paper. GPT-6.1 (OpenAI), working in Codex, recomputed the examples independently and wrote the interval computation of Section~\ref{sec:cert}. Both questions are due to Levin and Peres~\cite{LP}, the cycle case to Nacu~\cite{Nacu}, and the single-spin example on the same path to Shirai~\cite{Shirai}. The author directed the work and is responsible for the content of the paper and for decisions concerning its dissemination.

\section*{Data availability}
The Python scripts that produce every number in this paper, including the interval computation, are available at \url{https://github.com/yspennstate/ising-glauber-distance-monotonicity}.

\begin{thebibliography}{9}
\bibitem{LP} D.~A.~Levin and Y.~Peres, \emph{Markov Chains and Mixing Times}, second edition, with contributions by E.~L.~Wilmer, American Mathematical Society, Providence, RI, 2017.
\bibitem{Nacu} \c{S}.~Nacu, Glauber dynamics on the cycle is monotone, \emph{Probab. Theory Related Fields} \textbf{127} (2003), 177--185, \url{https://doi.org/10.1007/s00440-003-0279-x}.
\bibitem{Shirai} T.~Shirai, A remark on monotonicity for the Glauber dynamics on finite graphs, \emph{Proc. Japan Acad. Ser. A Math. Sci.} \textbf{86} (2010), 33--37, \url{https://doi.org/10.3792/pjaa.86.33}.
\bibitem{MathBrainRepo} Y.~Shmalo, \emph{Math Brain: public snapshot, code and tests}, 2026, \url{https://github.com/yspennstate/math-brain-public/tree/8c9a6bd478c4ec4082347151b7ae16db50c917dd}.
\bibitem{MathBrainPaper} Y.~Shmalo, \emph{Math Brain: shared, auditable memory for multi-agent mathematical research}, October 2026, \url{https://doi.org/10.5281/zenodo.24150837}.
\bibitem{BrainFramework} Y.~Shmalo, \emph{Brain Framework}, public software repository, 2026, \url{https://github.com/yspennstate/brain-framework-public/tree/5629e644cdd510173ae97e9aa92642dabd22dc57}.
\end{thebibliography}
\end{document}
